\chapter{Estado del arte}\label{cap:estado}

En esta sección se va a proceder a realizar un estudio del estado del arte en el ámbito de la planificación de 
recursos empresariales (Enterprise Resource Planning o ERP) para el sector de la oliva.Se analizarán cuáles 
son sus funcionalidades y limitaciones, así como las tecnologías utilizadas para su desarrollo.

\section{Contexto}

El desarrollo de software ERP para el sector de la oliva surge de la necesidad de integrar y controlar procesos 
muy específicos propios de una cadena agroindustrial compleja, caracterizada por una elevada regulación legislativa 
y por unos costes de producción y transformación significativos. La correcta gestión de la información se convierte 
en un factor clave para garantizar la eficiencia operativa, el cumplimiento normativo y la rentabilidad económica de 
las explotaciones, cooperativas y empresas del sector.

En este contexto, España se posiciona como líder mundial en la producción de aceite de oliva, concentrando de media en 
los últimos años más del 40\% de la producción mundial. A gran distancia le siguen otros países productores como Italia, 
con aproximadamente un 9,7\%, Grecia con un 8,2\% y Túnez con un 6,1\%, lo que pone de manifiesto la relevancia 
estratégica del sector oleícola español tanto a nivel económico como tecnológico \cite{principalesProductoresAceiteOliva}.
Esta posición de liderazgo implica, además, una mayor exigencia en términos de control, trazabilidad y calidad del 
producto final.

El proceso de producción de la oliva no puede considerarse lineal ni estándar, ya que intervienen múltiples fases 
interconectadas y con un alto grado de variabilidad. Desde la gestión de parcelas agrícolas y la recolección de la 
aceituna, pasando por la recepción, el almacenamiento en depósitos o bodegas, la molturación, el envasado y finalmente 
la comercialización, cada etapa genera información crítica que debe ser registrada y relacionada con el resto del 
proceso. Asimismo, en este sector conviven distintas entidades, como agricultores, cooperativas, almazaras, envasadoras
y comercializadoras, cada una con funciones específicas pero dependientes entre sí, lo que incrementa la complejidad de 
la gestión de los datos.

A esta complejidad productiva se suma una regulación especialmente estricta, orientada a garantizar la seguridad 
alimentaria, la trazabilidad y la calidad del aceite de oliva. Normativas como el Reglamento CE 178/2002 obligan a 
disponer de sistemas que permitan identificar el origen del producto y reconstruir su recorrido en cualquier punto de 
la cadena. Además, los consejos reguladores, las denominaciones de origen, los criterios de clasificación del aceite 
según su categoría y los análisis físico-químicos y organolépticos, así como las certificaciones de calidad y producción 
ecológica, imponen requisitos adicionales que deben ser gestionados de forma precisa y documentada.

Las soluciones ERP generalistas presentan importantes limitaciones a la hora de adaptarse a este contexto, ya que no 
contemplan de forma nativa conceptos fundamentales del sector oleícola como la gestión por campañas agrícolas, los 
rendimientos grasos, los lotes dinámicos, las mezclas de calidades, los trasiegos entre depósitos o el control de mermas.
Esto provoca que su implantación requiera un elevado nivel de personalización, incrementando los costes, los plazos y el 
riesgo de fracaso del proyecto.

Como consecuencia de estas necesidades específicas, se hace imprescindible el desarrollo de software ERP diseñado 
específicamente para el sector de la producción de oliva. (Estos sistemas se caracterizan por estar orientados a modelos 
de gestión por campañas, permitir un control exhaustivo de la trazabilidad desde la parcela hasta el producto final, 
gestionar lotes de forma dinámica y automatizar la generación de informes oficiales exigidos por las administraciones y 
organismos reguladores. Además, facilitan el control de costes, la gestión de la calidad y la integración de todos los 
agentes implicados en la cadena productiva, permitiendo a cooperativas, agricultores y empresas maximizar el valor de su 
producción y mejorar su competitividad en el mercado. QUITAR ??) A continuación, se va a proceder a analizar los softwares
que realizan esta labor para ver cómo resuelven la problemática mencionada anteriormente.

\section{Análisis de Softwares ERP para la producción del aceite de oliva}

Actualmente, podemos encontrar diversos softwares ERP que están directamente relacionados con la producción de oliva,
otros que no se relacionan directamente pero se ubican dentro del sector agrícola y otros que son genéricos pero se 
pueden adaptar al sector. Vamos a proceder a comentar algunos de ellos:

\begin{itemize}
    \item \textbf{Agroptima - Relación directa (https://www.agroptima.com)}
    \newline
    Agroptima es un Software de gestión agrícola cuyo objetivo es conseguir la simpleza a la hora de gestionar los cultivos
    y explotaciones por parte de agricultores y empresas agrícolas (según declara la propia empresa).

    En términos generales destaca por proporcionar un cuaderno de campo al agricultor, permitirle llevar un control de los 
    gastos y costes agrícolas (como los trabajadores) y saber en todo momento en qué estado se encuentra su finca.

    Entrando en más detalle, entre sus funcionalidades se encuentran:

    \begin{itemize}
        \item Localización de campos y cultivos, consulta de superficies y códigos SIGPAC.
        \item Anotación de toda la información que el agricultor considere relevante desde el móvil y sin cobertura.
        Entre esta información se destaca: trabajadores, maquinaria, horas trabajadas y productos.
        \item Toda actividad realizada y anotada quedará en el sistema y servirá para mejorar la toma de decisiones y 
        aumentar la rentabilidad de la gestión de la finca. 
        \item Las fincas o cultivos se podrán importar en la aplicación con un archivo Excel o Shapes, permitiendo visualizar
        las parcelas coloreadas según el tipo de cultivo y su superficie total.
        \item Gestión de roles, permitiendo que el usuario pueda gestionar aquellas actividades relacionadas con el.
    \end{itemize}

    Agroptima también destaca por ofrecer a sus clientes la posibilidad de generar un Cuaderno de Campo.

    \item \textbf{AgriCloud - Relación directa (https://www.agricloud.es)}
    \newline
    AgriCloud es un Software de gestión agrícola diseñado para llevar un control de las fincas, trabajos, personal y
    facturación, pensada para agricultores, cooperativas, empresas de servicios agrícolas y trabajadores por campaña.
    AgriCloud destaca por ofrecer funcionalidades como las siguientes:
    \begin{itemize}
        \item Un cuaderno de campo digital según la normativa (muy importante como hemos mencionado anteriormente).
        \item Control horario y partes de trabajo desde el móvil.
        \item Facturación electrónica conectada con Hacienda, permitiendo llevar un control del personal y ver la
        rentabilidad real por finca o cultivo.
        \item Costes ingresos e informes automáticos en un solo lugar.
        \item Acceso desde cualquier dispositivo, tanto para móvil como tablet u ordenador.
        \item Generación de un cuaderno de campo digital, facilitando la gestión y el seguimiento de las actividades
        agrícolas de forma eficiente y precisa.
    \end{itemize} 

    Como hemos comentado, AgriCloud está pensada para poder satisfacer las necesidades de varios sectores:
    \begin{itemize}
        \item Agricultores individuales: podrán registrar en cuestión de segundos sus tareas pendientes o realizadas,
        llevarán un control de costes e ingresos por cultivo (sabiendo cuánto cuesta y cuánto produce por cada parcela)
        y podrán generar informes de forma automática.
        \item Cooperativas y asociaciones: destaca por aportar funcionales como el registro y consulta de la información
        de cada miembro y su producción, supervisión en tiempo real de cada tarea y de los resultados de la campaña,
        generación de informes globales o por socio en PDF o Excel, permite la comunicación y coordinación entre el 
        personal a través de un entorno digital único.
        \item Empresas de servicios agrícolas: las empresas podrán gestionar partes digitales por cliente o finca, 
        convertir dichos partes en facturas electrónicas al instante, control de maquinaria y materiales (horas y costes
        asociados a cada servicio o cliente), ayuda en la toma decisión gracias a la consulta de resultados por tipo de 
        trabajo y campañas.
        \item Trabajadores por campaña: AgriCloud ofrece una gestión de las ``cuadrillas`` de trabajadores a través de un 
        fichaje y control del horario digital, asignación por tareas por finca o cultivo, accesos personalizados (cada perfil
        está adaptado al rol del empleado) y ofrece un seguimiento en la nube, consiguiendo que toda la información se 
        sincronice automáticamente.
    \end{itemize}

    \item \textbf{TuOlivar - Relación directa (https://www.tuolivar.com)}
    \newline
    TuOlivar es un Software de gestión agrícola para la producción del olivar. Se destacan principalmente por ofrecer las
    siguientes funcionalidades:
    \begin{itemize}
        \item Gestión de la explotación: permiten el control de fincas, parcelas, maquinaria, herramientas, mantenimiento
        de la maquinaria, trabajadores y campañas.
        \item Alta organización y uso desde cualquier ubicación: permite llevar un control de los trabajos realizados, control
        de gastos e ingresos y de la cosecha realizada, pudiendo agrupar la información por campañas, fincas y categorías.
        \item Gestión total: permite dar de alta a los clientes, almazaras, proveedores, notas, documentos o avisos, así como
        generar distintos informes para su control de costes y cosechas.
    \end{itemize}
    
    \item \textbf{Agroex (https://agroex.es)}
    \newline
    Agroex es un Software de gestión agrícola industrial y ganadería, enfocado en empresas grandes. Cuentan con un software
    adaptado a las necesidades específicas del sector, permitiendo optimizar recursos y ahorrar costes (tanto en campaña
    agrícola específica como a lo largo del año), gracias a la administración y gestión tanto de los cultivos como del 
    personal. Además, cuenta con integraciones como AEMET o MAGRAMA, permitiendo el acceso a datos en tiempo real que 
    permite ayudar en la toma de decisiones.
    \newline
    Agroex ofrece soluciones para distintos sectores, desde un enfoque en la explotación agrícola y consultoría agraria
    a otros como bodegas, viñedos y viveros.
    \newline
    Si analizamos las funcionalidades que ofrece este Software, Agroex destaca por:
    \begin{itemize}
        \item Administración y gestión de explotaciones agrícolas y agroindustriales: permite la gestión de cultivos y 
        administración de explotaciones a través de cuadernos de campo, gestión de plazos, riegos y recursos hídricos.
        \item Conectividad con entidades como MAGRAMA o AEMET para obtener información meteorológica precisa y en tiempo real.
        \item Gestión de costes.
        \item Gestión contable, generando de manera automática los apuntes contables para cada movimiento realizado.
        \item Gestión de almacén, controlando el stock y gestión de mercancías.
        \item Gestión documental, incluyendo la gestión de maquinaria, productos fitosanitarios y gestión de personal.
        \item Gestión económica y tesorería, generando informes completos y exactos de la liquidez y resultados económicos.
        \item Gestión de producción, a través del control de las actividades de fabricación, producción y manufactura.
        \item Gestión eficiente de los recursos hídricos.
        \item Recursos humanos, gracias a una conexión con la Seguridad Social.
        \item Facturación, ofreciendo herramientas para gestionar presupuestos, seguimiento comercial y factura electrónica.
    \end{itemize}

    \item \textbf{Galdón Software (Enfoque en Cooperativas) (https://www.galdon.com/erp-almazaras-envasadoras-aceite)}
    \newline
    Galdón Software es una empresa de ingeniería que ofrece soluciones software de gestión integrales para todo tipo de
    empresas e instituciones. Ofrecen servicios de consultoría y soluciones centralizadas para áreas de distintos
    negocios: ERP y CRM, Contabilidad y Finanzas, SGA, Apps, E-commerce y Recursos Humanos (RRHH).

    Dentro de los sectores en los que operan, desarrolladn Software ERP para almazaras y envasadoras de aceite. Con este
    ERP permiten llevar un control de la trazabilidad, producción y distribución del aceite y está pensado para cubrir
    las necesidades de cualquier empresa de aceite (oliva, girasol, soja, etc), almazara y envasadora.

    Entre sus características destacan:

    \begin{itemize}
        \item Planning de entrada y salida de mercancías.
        \item Control de necesidades de materias primas.
        \item Generación automática de partes de envasado y planificación asistida.
        \item Businnes Intelligence: cuentan con herramientas que permiten mostrar comparativas gráficas de ventas y 
        beneficios y análisis de ventas.
        \item Gesitón de cosecheros y puestos de compra.
    \end{itemize}

    Empresas reconocidas en el sector como Maeva (Granada) y Aceites Vallejo trabajan con este Software.

    \item \textbf{Campogest (https://www.campogest.com)}
    \newline
    Campogest es un Software diseñado por la empresa Hispatec y destinado para los Técnicos Agrícolas, permitiéndoles 
    realizar una gestión de sus actividades a través del móvil o tablet.

    Esta aplicación destaca por ofrecer las siguientes funcionalidades:

    \begin{itemize}
        \item Gestión de fincas, cultivos, previsión meteorológica, tratamientos fitosanitarios y recomendaciones personalizadas
        a través del móvil.
        \item Generación de un cuaderno de campo teniendo en cuenta las recomendaciones en materia de Normativa Legal y 
        Protocolos de Calidad.
        \item Generación de planes de abonado y notificaciones a través del correo electrónico al agricultor.
        \item Posibilidad de generar la finca o parcela a través de la cartografía SIGPAC, pudiendo asignar un propietario.
        Gracias a esta funcionalidad los agricultores pueden dar de alta sus fincas y cultivos desde el mapa.
    \end{itemize}

    \item \textbf{Visual NACERT (https://visualnacert.com)}
    \newline
    Visual Nacert es una empresa que diseña soluciones digitales de gestión para empresas del sector agroalimentario,
    agricultores e industria. Cuentan con una aplicación agrícola, a la que se puede acceder a través de móvil o 
    tablet, con o sin cobertura.

    Desde esta aplicación, permiten al agricultor/cliente realizar las funciones:

    \begin{itemize}
        \item Gestión de costes insumos, mano de obra y maquinaria, generando gráficos a partir de estos datos e 
        informes por parcela.
        \item Gestión del equipo de trabajo y control de stock.
        \item Gestión del Cuaderno de Campo Digital SIEX (de acuerdo a la normativa).
        \item Generación de informes de sostenibilidad.
    \end{itemize}

    Además de las funcionalidades anteriores, también ofrecen otras más centradas en la ayuda de toma de decisión al
    agricultor. Entre estas funcionalidades destacan:

    \begin{itemize}
        \item Informes de ayuda a la decisión basados en el análisis de datos del negocio y actividades realizadas 
        en campo.
        \item Análisis de la finca / parcela para ayudar en la toma de decisiones estratégicas sobre qué cultivar.
        \item Herramienta de IA que predice las necesidades de riego del cultivo y ofrece recomendaciones para una
        gestión del agua eficiente.
    \end{itemize}

    \item \textbf{AM System (https://amsystem.es)} 
    \newline
    AM System es una empresa que desarrolla software de gestión empresarial que ayuda a las empresas tanto en la 
    gestión como en la toma de decisiones. Dentro de las soluciones que ofrecen, cuentan con un software de gestión
    para las almazaras.

    Este software se basa en facilitar el trabajo diario en una almazara a través de la automatización de las tareas
    administrativas y de gestión. Entre sus características destacan: seguimiento del producto (entrada aceituna
    hasta su comercialización), gestión de fincas, control de depósitos y almacenes, liquidaciones y generación
    automática de cargos (entre otras).

    \item \textbf{BlueAiOlive}
    \newline 
    ``BlueAiOlive es un software de gestión para almazaras y envasadoras de aceite, diseñado para cubrir todas
    las necesidades de gestión del sector agrícola``. Así es como la empresa BlueAiCor define su ERP, destacando que a parte de cubrir
    las funcionalidades básicas de un ERP, tales como control y administración empresarial, también cuentan con módulos
    avanzados para la producción, envasado y comercialización del aceite de oliva.

    Para poder satisfacer dichas funcionalidades, cuentan con una herramienta integral que centraliza procesos, optimiza
    recursos y ofrece una visión global del negocio. De esta forma las empresas pueden tener un control total de cada etapa
    del proceso productivo.

    Entrando en el apartado del propio software, BlueAiOlive ofrece una app para oleicultores, donde se encontrará toda la
    información centralizada y disponible para el cliente a través de una interfaz sencilla e intuitiva. A través de ella,
    la empresa podrá llevar un control de la contabilidad, control del stock del almacén, facturación y liquidaciones, gestión
    de salidas de aceite, gestión de rendimiento de laboratorio (entre otras más).
\end{itemize}

\section{Comparativa con el estado del arte}

Rellenar con la información de qué ofrece mi producto en comparación con lo que se ha analizado que hay en el mercado

\section{Desarrollo de Software enfocado en Servicios Web}

Históricamente, el desarrollo de servicios web se fundamentó en la Arquitectura Orientada a Servicios (SOA). Como 
indicaban Ortiz y Riestra (2009) \cite{evolucionDinamicaSistemasSoftware}, el reto principal era lograr la 
interoperabilidad entre sistemas heterogéneos mediante protocolos robustos como SOAP y el uso de XML para el 
intercambio de datos. En aquel momento, el foco estaba en permitir que terminales con capacidades limitadas pudieran
acceder a lógica de negocio compleja residente en servidores.

No obstante, en la última década este paradigma ha evolucionado drásticamente. La rigidez de los protocolos basados 
en XML ha dado paso a la arquitectura REST (Representational State Transfer) y al formato JSON, que se han consolidado
como el estándar de facto por su ligereza y facilidad de consumo, especialmente en dispositivos móviles.

En la actualidad, han emergido nuevos paradigmas que complementan y, en muchos casos, sustituyen a las arquitecturas 
monolíticas tradicionales. Destacan especialmente los microservicios y las arquitecturas orientadas a eventos (
explicados a continuación), soluciones que han surgido como respuesta a la necesidad de construir sistemas con una 
mayor escalabilidad, flexibilidad y resiliencia ante grandes volúmenes de datos.

A continuación, se desglosarán los patrones de arquitectura software modernos, formatos de intercambio de datos y  
servicios web, ecosistemas de desarrollo web, desarrollo de aplicaciones móviles y por último, infraestructura
y seguridad. Este análisis permitirá examinar de manera detallada cómo han evolucionado los servicios web hasta 
alcanzar los estándares de eficiencia y robustez contemporáneos.

\subsection{Arquitecturas Software}

La arquitectura de software ha evolucionado progresivamente hasta convertirse en un ecosistema adaptable. El 
crecimiento exponencial en el desarrollo de servicios web y la necesidad de ofrecer soporte a aplicaciones móviles 
con alta disponibilidad han impulsado la aparición de nuevos enfoques y modelos arquitectónicos. En este contexto, 
la transición de sistemas centralizados (arquitecturas monolíticas) hacia sistemas distribuidos ha surgido de manera 
natural, motivada por la necesidad de desplegar funcionalidades de forma independiente y escalable.

De acuerdo con Richards y Ford \cite{libroArquitecturasSoftware}, las arquitecturas de software se pueden clasificar 
principalmente en dos categorías: monolíticas y distribuidas.

A continuación, se describen los estilos arquitectónicos más relevantes siguiendo la clasificación propuesta por los 
autores mencionados.

\subsubsection{Arquitectura Monolítica}

Es un tipo de arquitectura de software caracterizada principalmente por contener el código en una sola unidad de despliegue.
Dentro de este tipo de arquitectura se pueden encontrar los siguientes tipos:

\begin{itemize}
    \item \textbf{Arquitectura en capas}
    \newline
    Se utiliza como punto de partida ya que es el patrón más común y tradicional en el desarrollo de software. Este
    tipo de arquitectura se organiza en capas horizontales, donde cada una tiene un rol específico dentro de la 
    aplicación (por ejemplo, presentación y lógica de negocio). Aunque no hay un número fijo de capas, el estándar
    suele ser de cuatro:

    \begin{itemize}
        \item Capa de Presentación: representa la interfaz de usuario con la que se interacciona y es la encargada de
        manejar las peticiones que son realizadas al servidor.
        \item Capa de Negocio: contiene las reglas y la lógica del dominio de la aplicación.
        \item Capa de Persistencia: es la encargada de manejar los datos del sistema (interacción con la base de datos).
        \item Capa de Base de Datos: corresponde al motor de almacenamiento de datos físico, es decir, la propia base de datos.
    \end{itemize}

    Una de las características más relevantes de este tipo de arquitectura es que las capas están aisladas. Esto 
    significa que un cambio en una capa no debería afectar a las demás, siempre que la interfaz (el contrato) entre 
    ellas se mantenga. En este punto, se puede diferenciar entre capas de tipo ``cerrada``, que sería aquella que solo
    se puede comunicar con una inmediatamente inferior (tipo de capa por defecto), o de tipo ``abierta'', la cual permite
    que capas superiores puedan saltársela para acceder directamente a capas inferiores.

    Generalmente, este estilo se implementa como una única unidad de despliegue (monolito), lo que simplifica la gestión 
    inicial pero dificulta la escalabilidad individual de componentes.

    Valorando las características arquitectónicas se destaca:

    \begin{itemize}
        \item Costo y simplicidad: es una arquitectura muy fácil de entender y barata de implementar inicialmente.
        \item Facilidad de desarrollo: al ser un estándar tan conocido, los equipos pueden empezar a trabajar rápidamente.
        \item Escalabilidad y Elasticidad: esta características es, quizás, su mayor debilidad y esto ocurre porque al ser
        monolítica, no va a permitir un escalado de las partes específicas del sistema de forma independiente.
        \item Tolerancia a fallos: esta es otra debilidad que presenta, ya que si se produce un fallo en una de las capas
        (por ejemplo, base de datos), el fallo va a afectar a toda la aplicación.
    \end{itemize}

    Como conclusión, se puede destacar que esta arquitectura es ideal para aplicaciones pequeñas y sencillas donde el 
    presupuesto es limitado y la complejidad no justifica sistemas distribuidos más costosos.

    \item \textbf{Arquitectura Pipeline}
    \newline
    Este estilo de arquitectura se basa en el principio de computación de una secuencia de transformaciones. Su estructura
    es lineal y se compone de dos elementos principales: \textit{pipes} (tuberías) y \textit{filters} (filtros). Por un lado,
    los \textit{pipes} representan los canales de comunicación entre las etapas de procesamiento, suelen ser unidireccionales 
    y su función es puramente el transporte de datos de un filtro a otro; mientras que por el otro lado, los \textit{filters}
    son componentes independientes que realizan una tarea específica de procesamiento o transformación de los datos que reciben 
    del \textit{pipe}. Además, éstos son autónomos y no deben conocer la existencia de otros filtros en la cadena.

    Continuando con los filtros, éstos se pueden clasificar en distintos tipos según su rol funcional dentro de la topología:

    \begin{itemize}
        \item \textit{Producer (source)}: es el punto de inicio que genera los datos o los lee.
        \item \textit{Transformer}: es el encargado de realizar transformaciones de datos.
        \item \textit{Tester}: valida o filtra los datos baśandose en criterios específicos.
        \item \textit{Consumer (Sink)}: destino final donde terminan los datos.
    \end{itemize}

    En cuanto a sus características, esta arquitectura destaca por un alto desacoplamiento, ya que al ser los filtros
    independientes, es fácil intercambiar, añadir o modificar etapas en el pipeline sin afectar al resto del sistema. 
    Además, este estilo es ``\textit{stateless}'' entre filtros, es decir, cada filtro procesa lo que recibe y lo pasa 
    al siguiente sin mantener un estado global complejo.

    Valorando sus características arquitectónicas más relevantes, se clasifican en:

    \begin{itemize}
        \item Costo y simplicidad: es un estilo relativamente económico y fácil de implementar para los casos de uso 
        adecuados.
        \item Facilidad de desarrollo: la separación clara de responsabilidades en filtros independientes facilita 
        mucho el trabajo de codificación.
        \item Escalabilidad: al ser generalmente una arquitectura monolítica, es difícil de escalar de forma elástica.
        \item Tolerancia a fallos: si un filtro o una tubería falla, todo el proceso suele detenerse, lo que lo hace 
        poco resiliente.
    \end{itemize}

    Como conclusión, esta arquitectura es ideal para aplicaciones que requieren un procesamiento de datos secuencial y 
    discreto, como herramientas de procesamiento de texto, compiladores o sistemas de extracción, transformación y 
    carga (ETL) sencillos.

    \item \textbf{Arquitectura Microkernel}
    \newline
    También conocida como \textit{plug-ins}, es un estilo de arquitectura diseñado para productos software que se pueden
    instalar (como navegadores o IDEs). Su estructura se divide en dos componentes principales: \textit{Core System} 
    (sistema central), el cual contiene toda la lógica mínima necesaria para que la aplicación funcione, encargándose de 
    la ejecución básica y la gestión de los plug-ins; y \textit{Plug-in Components}, que son módulos independientes entre sí
    con lógica adicional, características especiales o adaptadores personalizados para extender el sistema central.

    En cuanto a su funcionamiento, el sistema central necesita saber qué plug-ins están disponibles. Para ello, utiliza  
    un registro o catálogo que contiene información sobre el plug-in, cómo acceder a él y qué protocolos utiliza. Por
    parte del \textit{plug-in}, para que este funcione debe cumplir con un contrato (interfaz) estándar definido por 
    el sistema central. 

    Valorando sus características arquitectónicas más relevantes, se clasifican en:

    \begin{itemize}
        \item Costo y simplicidad: es relativamente sencillo de implementar y muy eficiente en términos de costo para 
        productos que requieren extensibilidad.
        \item Evolutividad: es su mayor fortaleza ya que permite añadir funcionalidades de forma aislada y rápida a través 
        de plug-ins.
        \item Facilidad de prueba (testability): los plug-ins se pueden probar de forma aislada del sistema central.
        \item Escalabilidad y tolerancia a fallos: depende de la implementación, pero al ser frecuentemente monolítico, 
        no escala tan bien como los sistemas microservicios.
    \end{itemize}

    Como conclusión, es una arquitectura ideal para aplicaciones que necesitan evolucionar constantemente añadiendo nuevas 
    funciones o lógica según el cliente o el caso de uso, sin tener que reescribir el núcleo del sistema.

\end{itemize}

\subsubsection{Arquitectura Distribuida}

En esta arquitectura, a diferencia de la monolítica, los componentes van a estar separados y se conectarán a través de 
protocolos de acceso remoto.
Entre sus tipos, se pueden destacar los siguientes:

\begin{itemize}
    \item \textbf{Arquitectura basada en servicios}
    \newline
    Este estilo se considera un híbrido que combina aspectos de la arquitectura monolítica y los microservicios, siendo 
    una opción extremadamente popular para aplicaciones empresariales. 

    A diferencia de los microservicios, este estilo se basa en la creación de servicios de dominio de granularidad gruesa, 
    refiriéndose a servicios más grandes que los microservicios y menos detallados, que encapsulan una funcionalidad de negocio 
    más amplia y compleja. En una aplicación típica, suele haber entre 4 y 12 servicios, siendo cada servicio una unidad de 
    despliegue independiente que contiene toda la lógica necesaria para un dominio de negocio específico. Sin embargo, la base 
    de datos suele ser centralizada y compartida con todos (aunque si es necesario se puede ``partir''), al igual que la interfaz 
    de usuario, que hay una y es la que se comunica con los servicios.

    Este estilo es calificado en el libro como uno de los más equilibrados, destacando:

    \begin{itemize}
        \item Agilidad y desplegabilidad: al tener servicios independientes, es fácil realizar cambios y desplegarlos 
        rápidamente.
        \item Facilidad de prueba: los servicios de dominio pueden probarse de forma aislada, lo que mejora la calidad 
        del software.
        \item Costo y simplicidad: es mucho menos complejo y costoso de implementar que una arquitectura de microservicios pura.
        \item Escalabilidad y elasticidad: aunque es mejor que una arquitectura de capas, la base de datos compartida suele actuar 
        como un cuello de botella que limita el escalado infinito.
        \item Tolerancia a fallos: el fallo de un servicio no detiene todo el sistema, pero un fallo en la base de datos compartida 
        uede ser crítico.
    \end{itemize}

    Como conclusión, este estilo es usado para aplicaciones que necesitan un buen equilibrio entre agilidad, costo y escalabilidad, 
    especialmente cuando el dominio del negocio tiene un acoplamiento semántico que dificultaría mucho el uso de microservicios.

    \item \textbf{Arquitectura diriga por eventos} 
    \newline
    A diferencia del modelo tradicional basado en peticiones (\textit{request-based}), donde un componente solicita datos de forma 
    determinista, este modelo reacciona a situaciones que ocurren en el sistema (eventos). Se compone de componentes de procesamiento 
    de eventos desacoplados que operan de manera asíncrona. 

    Se van a distinguir dos formas fundamentales de implementar este estilo:

    \begin{itemize}
        \item Topología de Broker (intermediario)
        \newline
        Esta topología está caracterizada por no tener un orquestador central, siendo el flujo de mensajes distrbuido
        gracias a un \textit{broker} ligero. Funciona mediante una cadena de eventos: un componente procesa un evento 
        iniciador, publica lo que hizo como un "evento de procesamiento", y otros componentes interesados reaccionan a él.

        \item Topología de Mediador
        \newline
        Esta topología se utiliza cuando se requiere un control sobre el flujo de trabajo. Para ello, se cuenta con un
        ``Mediador de Eventos'' que será el encargado, conociendo toda la lógica de negocio, de coordinar los pasos 
        necesarios para procesar un evento inicial.        
    \end{itemize}

    Este estilo está caracterizado por ser asíncrono, permitiendo de esta manera que el sistema responda al usuario casi 
    instantáneamente mientras el procesamiento pesado ocurre en segundo plano. Esto evita bloqueos pero requiere del uso
    de técnicas (como colas de errores) que permitan gestionar fallos sin detener el sistema y que garanticen que los mensajes
    no se pierdan.

    Según el libro, la calificación general que se da para este estilo es la siguiente:

    \begin{itemize}
        \item Escalabilidad y rendimiento: son sus mayores fortaleza, ya que al ser asíncrona y desacoplada permite procesar 
        volúmenes masivos de datos.
        \item Agilidad y desplegabilidad: al estar desacoplados, los componentes pueden cambiar y desplegarse con facilidad.
        \item Simplicidad y facilidad de prueba: es un estilo muy complejo de diseñar, implementar y, sobre todo, de probar 
        debido a su naturaleza no determinista y distribuida.
    \end{itemize}

    Como conclusión, su uso es recomendado en aplicaciones que necesitan reaccionar en tiempo real a grandes flujos de información, 
    como en sistemas de procesamiento de datos de telemetría.

    \item \textbf{Arquitectura de microservicios}
    \newline

    Richards y Ford la definen como la "estrella brillante" del ecosistema actual, diseñada específicamente para abordar los problemas 
    de escalabilidad y agilidad de los sistemas monolíticos.

    Esta arquitectura se basa en el concepto de desacoplamiento extremo, tanto a nivel técnico como de dominio. A diferencia de otros 
    estilos, su objetivo principal no es la reutilización (que a menudo se evita para no crear acoplamiento), sino la separación de 
    intereses para permitir cambios rápidos. Aquí aparece una entidad clave que es el ``microservicio'', siendo una unidad física 
    separada que se puede desplegar, escalar y actualizar sin afectar a los demás, y cuya funcionalidad será única.

    Para garantizar el desacoplamiento de este estilo, cada microservicio posee sus propios datos. De esta manera, ningún otro servicio 
    puede acceder directamente a la base de datos de otro; la comunicación debe ser a través de interfaces (APIs).

    Entre sus componentes destacan:

    \begin{itemize}
        \item \textit{API Gateway}: actúa como punto de entrada para los clientes, manejando tareas como métricas, seguridad y enrutamiento.
        \item \textit{Sidecars}: componentes adjuntos a los servicios que manejan tareas transversales (comunicación, monitoreo) sin 
        afectar a la lógica de negocio.
        \item \textit{Service Mesh}: infraestructura para gestionar la comunicación entre todos los microservicios de forma consistente.
    \end{itemize}

    Aunque parezcan todo puntos a favor en esta arquitectura, es uno de los estilos más difíciles de implementar debido a la alta necesidad
    de gestión de los datos entre múltiples servicios y la complejidad de verificar que todo el ecosistema funcione correctamente. Además,
    requiere de conocimientos en áreas como Devops, automatización y monitoreo.

    Según el libro, su calificación es la siguiente:

    \begin{itemize}
        \item Escalabilidad y elasticidad: es el mejor estilo para sistemas que necesitan crecer de forma masiva y dinámica.
        \item Agilidad y desplegabilidad: Permite ciclos de entrega continuos.
        \item Tolerancia a fallos: el aislamiento asegura que el fallo de un servicio no derribe todo el sistema (siempre 
        que se implementen patrones de resiliencia).
        \item Costo y simplicidad: es el estilo más caro y complejo de desarrollar, mantener y operar. No se recomienda 
        para aplicaciones sencillas.
    \end{itemize}

    Como solución, la arquitectura de microservicios es muy útil para resolver problemas de gran escala, pero conlleva un "impuesto de 
    complejidad" que solo debe pagarse si las necesidades de negocio lo justifican.
\end{itemize}

\subsection{Comunicación y Servicios Web}

La comunicación entre componentes es el elemento que define la eficiencia de una arquitectura distribuida. Mientras que en el pasado 
la interoperabilidad se buscaba mediante protocolos estrictos basados en estándares corporativos, el auge de la web móvil ha impuesto 
la necesidad de protocolos de baja latencia y formatos de datos ligeros. Según Richards y Ford \cite{libroArquitecturasSoftware}, el 
éxito de una arquitectura distribuida depende de la gestión de los contratos de comunicación y el ancho de banda disponible.

Esta interacción es posible gracias a las Interfaces de Programación de Aplicaciones (\hyperlink{api}{API}). Entre sus tipos principales 
destacan:

\begin{itemize}
    \item \textbf{REST \textit{(Representational State Transfer)}}
    \newline
    Consolidado como el estándar de facto para servicios web. REST utiliza los métodos del protocolo HTTP para definir acciones sobre 
    los recursos del sistema a través de \textit{endpoints} (URLs).

    Los principales verbos HTTP y su propósito son: \cite{metodosHttp}

    \begin{itemize}
        \item \textit{GET}: solicita una representación de un recurso específico. Las peticiones que usan el método GET 
        sólo deben recuperar datos. 
        \item \textit{PUT}: reemplaza todas las representaciones actuales del recurso de destino con la carga útil de 
        la petición.
        \item \textit{POST}: se utiliza para enviar una entidad a un recurso en específico, causando a menudo un cambio 
        en el estado o efectos secundarios en el servidor.
        \item \textit{DELETE}: borra un recurso en específico.
        \item \textit{HEAD}: pide una respuesta idéntica a la de una petición GET, pero sin el cuerpo de la respuesta.
        \item \textit{PATCH}: aplica modificaciones parciales a un recurso.
        \item \textit{OPTIONS}: utilizado para describir las opciones de comunicación para el recurso de destino.
        \item \textit{TRACE}: realiza una prueba de bucle de retorno de mensaje a lo largo de la ruta al recurso de destino.
        \item \textit{CONNECT}: establece un túnel hacia el servidor identificado por el recurso.
    \end{itemize}

    Para confirmar el resultado de estas acciones, el servidor devuelve códigos de estado HTTP, tales como: \texttt{200 OK} (éxito), 
    \texttt{400 Bad Request} (error del cliente) o \texttt{500 Internal Server Error} (error del servidor).

    \item \textbf{GraphQL} \cite{graphql}
    \newline
    GraphQL utiliza una sintaxis intuitiva que permite a los clientes enviar una única consulta GraphQL a una API y 
    recibir exactamente los datos necesarios (en lugar de acceder a endpoints complejos con muchos parámetros). 
    Esta obtención de datos más eficiente puede mejorar el rendimiento del sistema y la facilidad de uso para los 
    desarrolladores. 

    Es especialmente útil para crear APIs en entornos complejos con requisitos de interfaz que cambian rápidamente. Ni 
    las API de REST ni las de GraphQL son inherentemente superiores; son herramientas diferentes que se adaptan a 
    diferentes tareas.

    \item \textbf{gRPC} \cite{grpc}
    \newline
    Es un marco de comunicación de alto rendimiento desarrollado por Google. Utiliza como sistema de comunicación
    el protocolo RPC (\textit{Remote Procedure Call}) en lugar de texto, lo que lo hace significativamente más rápido 
    que REST para la comunicación interna entre microservicios.

    \item \textbf{SOAP} \cite{soap}
    \newline
    SOAP (anteriormente conocido como Simple Object Access Protocol) es un protocolo ligero para el intercambio de 
    información en entornos descentralizados y distribuidos. Es un protocolo basado en XML que define tres partes 
    en todos los mensajes: \textit{sobre}, que define una infraestructura para describir qué hay en un mensaje y 
    cómo procesarlo; \textit{reglas de codificación}, que expresa instancias de tipos de datos definidos por la 
    aplicación, y \textit{estilos de comunicación}, que pueden seguir el formato de llamada a procedimiento remoto 
    (RPC) o el formato orientado a mensajes (Documento). 
\end{itemize}

En cuanto al formato de intercambio de datos, Se ha producido una transición del uso de XML (extenso y costoso de 
procesar) al dominio absoluto de JSON (JavaScript Object Notation). JSON es un formato de texto ligero, fácil de 
leer para humanos y sumamente rápido de procesar por los motores modernos.

Finalmente, la comunicación entre sistemas va más allá del modelo tradicional de petición-respuesta de HTTP, existiendo 
otros protocolos para servicios que requieren inmediatez como \textbf{WebSockets}, proporcionando un canal de comunicación 
bidireccional constante entre el cliente y el servidor, y \textbf{Server-Sent Events}, permitiendo que el servidor envíe 
actualizaciones de forma automática al cliente sin que este tenga que solicitarlas repetidamente.

\subsection{Ecosistema del Desarrollo Web}

En este apartado se describen las principales tecnologías consideradas para el desarrollo del proyecto Olivaris. 
Estas se clasifican en dos grandes áreas: el Frontend (lado del cliente) y el Backend (lado del servidor).

\subsubsection{Tecnologías para el frontend}

El frontend comprende todo aquello con lo que el usuario interactúa directamente. En la última década, este campo ha 
evolucionado desde documentos estáticos hacia \hyperlink{singlePageApplication}{\textit{Single Page Application}} y 
\hyperlink{progressiveWebApps}{\textit{Progressive Web Apps}} de gran complejidad. El objetivo actual es gestionar 
estados complejos, asegurar la accesibilidad y optimizar el rendimiento bajo el paradigma de \hyperlink{clientSideRendering}{\textit{Client-Side Rendering}}  
(CSR) o enfoques híbridos.

Aunque hoy en día se utilizan muchas herramientas, la base de cualquier interfaz web descansa sobre tres pilares:

\begin{itemize}
    \item \textbf{HTML5 (\textit{HyperText Markup Language})} \cite{html}
    \newline
    Es el componente más básico de la Web. Define la estructura y el significado del contenido mediante etiquetas 
    (marcas), permitiendo que el navegador interprete textos, imágenes y multimedia.

    \item \textbf{CSS3 (\textit{Cascading Style Sheets})} \cite{css}
    \newline
    Lenguaje de estilos encargado de la presentación visual. Define cómo se renderizan los elementos en pantalla 
    (diseño, colores, tipografías y adaptabilidad).

    \item \textbf{JavaScript} \cite{javascript}
    \newline
    Si bien es más conocido como un lenguaje de \textit{scripting} para páginas web, y es usado en muchos entornos 
    fuera del navegador, JavaScript es un lenguaje de programación basada en prototipos, multiparadigma, de un solo 
    hilo, dinámico, con soporte para programación orientada a objetos, imperativa y declarativa. Gracias a JavaScript, se
    ha conseguido que la web se convierta en una pieza interactiva y dinámica para el usuario.

    Mención especial merece \textbf{TypeScript}, como un superset de JavaScript que añade tipado estático. Su uso 
    mejora la detección de errores en tiempo de desarrollo y facilita el mantenimiento de bases de código extensas.
\end{itemize}

El desarrollo ha evolucionado del uso de "Vanilla JS" (JS puro) hacia marcos de trabajo (\textit{\hyperlink{marcosTrabajo}{frameworks}}) 
que aumentan la eficiencia:

\begin{itemize}
    \item \textbf{React} \cite{react}
    \newline
    Librería basada en componentes reutilizables. Destaca por su flexibilidad y su capacidad para manejar interfaces 
    de usuario altamente dinámicas mediante un DOM virtual.

    \item \textbf{Angular} \cite{angular}
    \newline
    Angular es un framework integral que utiliza TypeScript de forma nativa. Su arquitectura modular y su sistema de 
    inyección de dependencias lo hacen ideal para aplicaciones empresariales escalables.
\end{itemize}

Para el ecosistema móvil, se han evaluado diversas estrategias que permiten alcanzar a usuarios de Android e iOS:

\begin{itemize}
    \item \textbf{Flutter} \cite{flutter}
    \newline
    Flutter es un framework de código abierto de Google para crear aplicaciones multiplataforma compiladas de forma 
    nativa a partir de una única base de código, ofreciendo un alto rendimiento y una interfaz consistente.

    \item \textbf{React Native} \cite{reactNative}
    \newline
    React Native es una librería de JavaScript que permite construir aplicaciones nativas para Android e iOS. Proporciona 
    componentes nativos que se asignan directamente a los bloques de construcción de la interfaz.

    \item \textbf{Kotlin} \cite{kotlin}
    \newline
    Kotlin es un lenguaje de programación de tipo estático utilizado para desarrollar aplicaciones móviles en Android,
    destacando por su seguridad y concisión.

    \item \textbf{Ionic} \cite{ionic}
    \newline
    Ionic es un SDK (kit de desarrollo de software) de código abierto que permite crear aplicaciones multiplataforma 
    utilizando una única base de código. La aplicación se construye una sola vez con tecnologías web (como HTML, CSS 
    y JavaScript), y se despliega en cualquier dispositivo. 
\end{itemize}

\subsubsection{Tecnologías para el backend}

El backend representa la capa de lógica de procesamiento y acceso a datos que reside en el servidor. El enfoque actual 
prioriza la escalabilidad, la capacidad de respuesta ante grandes volúmenes de datos y la exposición de servicios 
mediante APIs (\textit{Application Programming Interfaces}) robustas, generalmente bajo el estándar REST o GraphQL.

Dentro de los lenguajes más utilizados para el desarrollo del backend, se encuentran: 

\begin{itemize}
    \item \textbf{Java} 
    \newline
    Java se mantiene como un pilar en sistemas empresariales críticos debido a su robustez, tipado fuerte y excelente 
    gestión de memoria a través de la JVM (\textit{Java Virtual Machine}).

    Uno de los frameworks de Java más utilizados es Spring Boot \cite{springBoot}, siendo el framework de código abierto 
    líder, que facilita el uso de marcos de trabajo basados en Java para crear microservicios y aplicaciones web autónomas.

    En cuanto a sus ventajas, proporciona configuración automática de librerías, servidores embebidos (como Tomcat) y un 
    ecosistema maduro para la creación de entornos de pruebas. Además, Aumenta drásticamente la productividad al eliminar 
    la configuración manual repetitiva.

    \item \textbf{Python}
    \newline
    Es el lenguaje con mayor crecimiento actual debido a su sintaxis limpia y su dominio en Inteligencia Artificial y Ciencia 
    de Datos, lo que facilita la integración de modelos de \textit{Machine Learning} en el backend.

    Python también puede ser útil para el desarrollo web, ofreciendo una amplia gama de frameworks como Djano, Flask o
    FastApi.

    En cuanto a Django \cite{django}, es un framework para desarrollo de aplicaciones "con todo incluido" diseñado para el 
    desarrollo rápido. Utiliza una arquitectura basada en módulos reutilizables, lo que reduce significativamente el tiempo 
    de salida al mercado (\textit{Time-to-Market}). Es la base de plataformas de gran escala como Instagram.

    Por otro lado, FastApi \cite{fastapi}, siendo una opción moderna y de alto rendimiento basada en el tipado de Python. 
    Destaca por la generación automática de documentación técnica (OpenAPI/Swagger) y su soporte nativo para programación 
    asíncrona, lo que lo hace ideal para APIs de baja latencia.

    \item \textbf{Node.js} \cite{nodeJs}
    \newline
    Node.js no es un lenguaje en sí, sino un entorno de ejecución que permite ejecutar JavaScript en el servidor. Su 
    arquitectura está basada en un modelo de E/S (entrada/salida) sin bloqueo y orientado a eventos. Es extremadamente 
    eficiente para aplicaciones en tiempo real y servicios que manejan una gran cantidad de conexiones simultáneas, 
    ofreciendo una alta escalabilidad horizontal.

    Para su desarrollo web, cuenta con frameworks como \textbf{Express}, siendo un framework minimalista y flexible que 
    ofrece libertad total al desarrollador, \textbf{NestJs}, que destaca por ser progresivo y modular (inspirado en Angular) 
    que utiliza TypeScript por defecto, ideal para construir aplicaciones backend escalables y fáciles de mantener.
\end{itemize}

\subsection{Sistemas de almacenamiento}

El almacenamiento de datos ha evolucionado desde sistemas de archivos planos hacia Sistemas de Gestión de Bases de 
Datos (SGBD). En el desarrollo moderno, la elección de una base de datos no depende solo del volumen de información, 
sino de su estructura, la velocidad de consulta y el nivel de consistencia requerido. El estado del arte actual se 
divide principalmente en los modelos Relacionales (SQL) y No Relacionales (NoSQL).

\subsubsection{Bases de Datos Relacionales (SQL) \cite{baseDatosSQL}} 

Este modelo organiza los datos en tablas con relaciones predefinidas. Cada tabla representa una entidad, estructurando 
la información en filas (registros) y columnas (atributos). Su principal fortaleza es la integridad referencial y el 
uso del lenguaje estandarizado SQL (\textit{Structured Query Language}) para la manipulación de datos.

Estas bases de datos destacan por tener un esquema rígido (es necesario definir la estructura antes de insertar datos) 
y cumplimiento de las propiedades ACID (Atomicidad, Consistencia, Aislamiento y Durabilidad).

Entre las más destacadas se encuentran: MySQL / MariaDB, PostgreSQL y Microsoft SQL Server.

\subsubsection{Base de Datos NoSQL \cite{baseDatosNoSQL}}

Diseñadas para modelos de datos específicos, las bases de datos NoSQL ofrecen esquemas flexibles que permiten una alta 
escalabilidad horizontal. Son ideales para aplicaciones modernas que requieren un desarrollo rápido e iterativo.

Se clasifican según su forma de almacenar la información:

\begin{itemize}
    \item Orientadas a documentos: almacenan datos en formatos similares a JSON (como MongoDB).
    \item Clave - Valor: optimizadas para lecturas muy rápidas mediante una clave única (Redis).
    \item Orientadas a grafos: enfocadas en la relación entre nodos (como Neo4j).
\end{itemize}

Están caracterizadas por ofrecer escalabilidad horizontal mediante clústeres distribuidos y manejo 
eficiente de grandes volúmenes de datos no estructurados.

\subsubsection{Bases de Datos en la Nube \cite{baseDatosNube}}

Representan la evolución del almacenamiento hacia el modelo \textit{as-a-service}. Aunque pueden ser relacionales o NoSQL, 
se diferencian de las locales por aprovechar los beneficios de la computación en la nube:

\begin{itemize}
    \item Escalabilidad y Agilidad: permiten aumentar los recursos de forma elástica según la demanda.
    \item Reducción de costes: eliminan la necesidad de gestionar hardware físico y mantenimiento de infraestructura.
    \item Alta disponibilidad: ofrecen réplicas automáticas y copias de seguridad gestionadas por el proveedor (como AWS 
    RDS o Azure SQL).
\end{itemize}

\subsection{Despliegue y estrategias}

El despliegue de software (\textit{deployment}) es el proceso crítico que transforma el código fuente en una aplicación funcional 
y accesible para los usuarios finales. En la actualidad, se concibe como un ciclo continuo y automatizado que garantiza 
la disponibilidad y estabilidad del sistema. Este proceso se compone de las siguientes etapas fundamentales:

\begin{itemize}
    \item \textbf{Integración Continua (CI):} es la fase donde el código se compila, valida y somete a pruebas automáticas cada 
    vez que un desarrollador realiza un cambio, asegurando la integridad de la base de código.
    \item \textbf{Entrega/Despliegue continuo (CD):} una vez validado el código, se realiza el envío automático a los entornos 
    de prueba o producción.
    \item \textbf{Aprovisionamiento de infraestructura:} consiste en la creación y configuración de los recursos (servidores, 
    bases de datos, redes) donde residirá la aplicación, idealmente mediante Infraestructura como Código (IaC).
    \item \textbf{Monitorización y \textit{feedback}:} una vez en producción, la aplicación es vigilada para detectar errores en 
    tiempo real, permitiendo un \textit{rollback} (vuelta atrás) inmediato si se detectan anomalías.
\end{itemize}

Un aspecto fundamental es decidir que estrategia de despliegue se va a llevar a cabo cuando se realiza una aplicación 
y cómo se libera a los usuarios. Los tipos de estrategia pueden variar según las necesidades de la infraestructura y 
según la estrategia que se quiera seguir. 

Según la infraestructura, el despliegue se puede realizar en servidores físicos propios (\textit{on-premise}), a través 
de proveedores en la nube (como \textit{AWS}), o ejecutándose a través de eventos, sin gestión directa de servidores 
(\textit{serverless}). Por otro lado, según la estrategia de liberación que se siga, se encuentran los siguientes 
tipos: \cite{tiposDespliegue} 

\begin{itemize}
    \item \textit{Rolling upgrade}
    \newline
    Se realiza una actualización secuencial del servicio, actualizando uno a uno los componentes. Este tipo de despliegue 
    tiende a ser lento tanto en su ejecución como en su capacidad de volver atrás (\textit{rollback}) y generalmente se emplea en 
    combinación con otras estrategias.

    \item \textit{Blue/Green}
    \newline
    Estrategia muy común que implica el uso de dos versiones del servicio: la versión nueva (\textit{blue}) y la versión estable 
    existente (\textit{green}). Los usuarios continúan utilizando la versión verde mientras se prueba y evalúa la versión azul. 
    Cuando se considera lista, los usuarios cambian a la versión azul. Si se detecta algún problema, es posible regresar a la 
    versión verde.

    \item \textit{ Red/Black}
    \newline
    Similar a la estrategia blue/green, con la diferencia de que se desvía una pequeña porción del tráfico desde la 
    versión negra (estable) hacia la nueva versión. Luego de evaluar su rendimiento, se realiza un desvío aún más grande 
    del tráfico y, finalmente, se redirecciona todo el tráfico hacia la nueva versión. Se mantiene la versión negra en 
    línea para permitir un \textit{rollback} rápido si es necesario.

    \item \textit{Dark Launcher}
    \newline
    Se implementa una nueva versión en paralelo a la versión estable existente. La particularidad aquí es que se duplica 
    una parte del tráfico para observar cómo se comporta la nueva versión bajo una carga duplicada. Luego, se realiza un 
    redireccionamiento parcial del tráfico hacia la nueva versión.

    \item \textit{Canary Release}
    \newline
    Esta estrategia implica un despliegue parcial basado en el tiempo y el rendimiento de la nueva versión, consiguiendo
    de esta manera una implementación gradual y controlada.
\end{itemize}

En cuanto a las tecnologías y herramientas que se pueden utilizar para realizar el despliegue y gestionar una aplicación,
se pueden clasificar en:

\begin{itemize}
    \item Herramientas de contenerización, siendo las más utilizadas son Dokcer y Kubernetes. Por un lado, Docker 
    \cite{docker} empaqueta software en unidades estandarizadas llamadas contenedores que incluyen todo lo necesario 
    para que el software se ejecute, incluidas bibliotecas y herramientas de sistema, pudiendo ejecutar aplicaciones 
    en cualquier entorno. Por otro lado, Kuberntes \cite{kubernetes} facilita la automatización y la configuración 
    declarativa, además de administrar contenedores, orquestando la infraestructura de cómputo, redes y almacenamiento 
    para que las cargas de trabajo de los usuarios no tengan que hacerlo.
    \item Herramientas de automatización, como GitHub Actions o Jenkins que ejecutan los flujos de trabajo de forma automática.
    \item Herramientas que permiten la definición de la infraestructura como código (IaC), como Terraform, que permiten 
    definir y gestionar la infraestructura mediante archivos de configuración, facilitando la replicabilidad.
\end{itemize}

\subsection{Autenticación y autorización}

La seguridad es un apartado crítico en el desarrollo de cualquier sistema; es una responsabilidad transversal que se 
apoya en protocolos estandarizados para garantizar la interoperabilidad.

La autenticación (\textit{AuthN}) se define como el proceso técnico de verificar que un usuario es quien dice ser, 
constituyendo la primera barrera de defensa. A continuación, se detallan los mecanismos más utilizados:

\begin{itemize}
    \item \textbf{Basada en sesiones (formularios + cookies) \cite{authCookies}}
    \newline
    Es el método tradicional donde el usuario envía sus credenciales y, tras la validación, el servidor crea una sesión 
    y devuelve una cookie al navegador. En cada petición posterior, el cliente envía dicha cookie para que el servidor 
    valide el ID de sesión.

    Es un método que destaca su alta seguridad en aplicaciones web simples y capacidad de revocación instantánea del 
    acceso; sin embargo, presenta dificultades para escalar (debido al estado en el servidor).

    \item \textbf{Basada en tokens (JWT) \cite{authJWT}}
    \newline
    Es el estándar predominante en aplicaciones modernas y arquitecturas desacopladas. Tras la validación, el servidor 
    genera un JSON Web Token (\hyperlink{jwt}{JWT}) firmado digitalmente. El cliente lo almacena y lo incluye en la cabecera de cada 
    petición para acceder a recursos protegidos.

    Entre sus ventajas, destaca por una escalabilidad total (stateless) y compatibilidad nativa con aplicaciones móviles;
    sin embargo, presenta un riesgo de seguridad si el token es interceptado, ya que el acceso es válido hasta su 
    expiración (difícil de revocar antes de tiempo).

    \item \textbf{OAuth 2.1 y OpenID Connect \cite{oauth}}
    \newline
    OAuth 2.1 es un framework que permite el acceso delegado (una aplicación accede a datos de otra sin compartir contraseñas), 
    mientras que OpenID Connect (OIDC) es una capa de identidad sobre OAuth que habilita el \hyperlink{singleSignOn}{\textit{Single Sign-on}}.

    Este sistema de autenticación destaca por su gran escalabilidad e integración con múltiples proveedores (Google, GitHub, etc.);
    sin embargo, presenta una mayor complejidad técnica en su implementación inicial.

    \item \textbf{Multifactor \cite{multifactor}}
    \newline
    Se trata de un mecanismo que combina diversos factores de verificación: conocimiento (PIN/contraseña), posesión (dispositivo 
    móvil) e inherencia (biometría).

    Como ventaja destaca por mitigar riesgos de robo de credenciales o tokens; aunque un atacante obtenga el JWT, carecería del 
    segundo factor; sin embargo, Puede degradar la experiencia de usuario (UX) al añadir fricción y lentitud al proceso de entrada.
\end{itemize}

\subsection{Elección de tecnologías}

Tras el análisis de las opciones disponibles, se han seleccionado las siguientes tecnologías para el desarrollo del proyecto Olivaris, 
buscando un equilibrio entre rendimiento, seguridad y mantenibilidad.

\begin{itemize}
    \item \textbf{Frontend: React Native}
    \newline
    El acceso a la plataforma será principalmente a través de dispositivos móviles, dada la naturaleza del trabajo agrícola. La movilidad 
    y disponibilidad inmediata que ofrece un smartphone son claves para que los agricultores interactúen con el sistema a pie de campo.

    Las razones de la elección han sido:

    \begin{itemize}
        \item Utiliza JavaScript/TypeScript, uno de los lenguajes con mayor comunidad y demanda laboral, superando actualmente el ecosistema 
        de Flutter en oportunidades regionales.
        \item A diferencia de las soluciones híbridas, React Native utiliza componentes nativos reales de iOS y Android, respetando las guías 
        de diseño y ofreciendo una experiencia de usuario fluida.
        \item El uso de Expo facilita la configuración del entorno (evitando dependencias complejas de Android Studio) y permite actualizaciones 
        \textit{Over-The-Air} (OTA), permitiendo corregir errores sin que el usuario tenga que descargar una nueva versión de las tiendas oficiales.
        \item El ecosistema de librerías es muy amplio. Si se necesita integrar un SDK de terceros (pagos, mapas, analíticas), es muy probable que ya exista
        un paquete oficial o muy probado para React Native.
    \end{itemize}

    \item \textbf{Backend: Java con SpringBoot}
    \newline
    Para la lógica de servidor se ha optado por el ecosistema de Java, priorizando la robustez y la seguridad del sistema. Las razones por las que se ha elegido
    son las siguientes:

    \begin{itemize}
        \item El tipado estático de Java permite detectar la mayoría de los errores en tiempo de compilación. Esto ofrece una seguridad que lenguajes dinámicos 
        como Python o JavaScript no pueden garantizar en la misma medida.
        \item Java maneja hilos de ejecución de forma eficiente, siendo superior en tareas que requieren un uso intensivo de CPU o procesamiento paralelo.
        \item Spring Boot ofrece una estructura de proyecto coherente y estandarizada. La gestión de acceso a datos y seguridad se realiza bajo patrones conocidos, 
        facilitando que cualquier desarrollador se integre rápidamente en cualquier proyecto.
        \item Cuenta con \textit{Spring Security}, uno de los marcos de seguridad más avanzados y granulares del mercado.
        \item El sistema está preparado para manejar transacciones ACID complejas. Además, la facilidad de Java para generar reportes en PDF y su compatibilidad 
        con sensores IoT aseguran que el proyecto pueda crecer sin limitaciones técnicas.
    \end{itemize}

    Teniendo en cuenta lo anterior, aunque inicialmente la carga de procesamiento no sea muy alta, SpringBoot permite manejar las transacciones ACID de forma
    impecable y este proyecto, con una sola acción se van a desencadenar diferentes acciones ``por debajo''. Además, Java cuenta con librerías muy potentes
    para generar PDFs y si en un futuro se deciden añadir funcionalidades extras, la escalabilidad es muy buena (integración con sensore IoT).

    También lo he decidido así porque Java con SpringBoot es una tecnología muy demandada en el mercado laboral actualmente, por lo que considero que es un buen
    punto para empezar a aprenderlo.

    \item \textbf{Base de datos: PostgreSQL}
    \newline
    Se ha optado por un modelo relacional (SQL) debido a la complejidad de las asociaciones y dependencias jerárquicas entre las entidades del dominio de Olivaris.

    El uso de PostgreSQL en lugar de MySQL o MariaDB se debe a las siguientes razones:

    \begin{itemize}
        \item Existen relaciones jerárquicas, por lo que un modelo relacional es el más eficiente para garantizar la integridad de los datos cuando existen múltiples 
        tablas interconectadas.
        \item PostgreSQL destaca sobre otros motores SQL por su soporte avanzado de tipos no estructurados (\hyperlink{jsonb}{JSONB}), permitiendo lo mejor de los dos 
        mundos: la rigidez del SQL y la flexibilidad de NoSQL en campos específicos.
        \item Su estricto cumplimiento de los principios ACID (Atomicidad, Consistencia, Aislamiento y Durabilidad) garantiza que la información nunca se corrompa, 
        incluso ante fallos del sistema.
    \end{itemize}
\end{itemize}
